% บทที่ 2 เอกสารและงานวิจัยที่เกี่ยวข้อง
% ดึงจาก Google Doc ของทีม (Collab_Reflect_Report) เฉพาะบทที่ 2 เมื่อ 2026-11-05
% รายการอ้างอิงอยู่ใน sampleReport.bib
\chapter{\ifenglish Literature Review\else เอกสารและงานวิจัยที่เกี่ยวข้อง\fi}\label{ch:background}

โครงงาน Collab Reflection เป็นเว็บแอปพลิเคชันสำหรับสนับสนุนการสะท้อนคิดและการประเมินตนเองและเพื่อนร่วมทีม (Self- and Peer Assessment) โดยมุ่งเน้นการพัฒนาทักษะการทำงานเป็นทีมของนักศึกษาในรายวิชาที่มีผู้เรียนจำนวนมาก บทนี้นำเสนอแนวคิดและงานวิจัยที่ใช้เป็นพื้นฐานในการกำหนดแนวทางออกแบบระบบ ได้แก่ การเรียนรู้แบบร่วมมือ ปัญหาที่เกิดขึ้นในการทำงานกลุ่ม การประเมินโดยเพื่อนและการสะท้อนคิด เกณฑ์ประเมิน CATME แนวทางจัดการข้อเสนอแนะที่อ่อนไหว การคุ้มครองข้อมูลส่วนบุคคล ตลอดจนเทคโนโลยีและสถาปัตยกรรมที่เกี่ยวข้องกับการพัฒนาระบบ

\section{การทำงานเป็นทีมและการเรียนรู้แบบร่วมมือ}

การเรียนรู้แบบร่วมมือ (Cooperative Learning) เป็นการจัดการเรียนรู้ที่ให้ผู้เรียนทำงานร่วมกันเป็นกลุ่มขนาดเล็กเพื่อบรรลุเป้าหมายร่วมกัน สมาชิกแต่ละคนต้องพึ่งพาและสนับสนุนกัน โดยมีองค์ประกอบสำคัญ ได้แก่ การพึ่งพากันในเชิงบวก (positive interdependence) ความรับผิดชอบรายบุคคล (individual accountability) การปฏิสัมพันธ์ระหว่างสมาชิก ทักษะทางสังคม และการประมวลผลกระบวนการกลุ่ม (group processing)~\cite{johnson1998active} องค์ประกอบเหล่านี้ชี้ให้เห็นว่าความสำเร็จของการทำงานกลุ่มไม่ได้พิจารณาจากผลงานสุดท้ายเพียงอย่างเดียว แต่ควรพิจารณาพฤติกรรมและการมีส่วนร่วมของสมาชิกแต่ละคนด้วย

สำหรับรายวิชาที่มีการทำงานกลุ่ม ผู้สอนอาจเห็นผลงานที่ส่งได้ชัดเจนกว่ากระบวนการทำงานระหว่างสมาชิก โดยเฉพาะกิจกรรมที่เกิดขึ้นนอกชั้นเรียน การมีเครื่องมือให้สมาชิกประเมินตนเองและสะท้อนพฤติกรรมของเพื่อนจึงช่วยเพิ่มข้อมูลประกอบการติดตามการทำงานเป็นทีม ซึ่งเป็นแนวคิดหลักของ Collab Reflection

\subsection{ขั้นตอนพัฒนาการของกลุ่ม}

Tuckman เสนอว่ากลุ่มมักพัฒนาผ่าน 4 ขั้นตอน ได้แก่ Forming (การเริ่มรวมกลุ่ม) Storming (การเผชิญความเห็นต่าง) Norming (การสร้างบรรทัดฐานร่วมกัน) และ Performing (การทำงานเพื่อบรรลุเป้าหมาย)~\cite{tuckman1965developmental} ในช่วงเริ่มต้น สมาชิกอาจยังไม่เข้าใจบทบาทและความคาดหวังของกันและกัน ส่วนช่วง Storming อาจเกิดความขัดแย้งเกี่ยวกับหน้าที่ วิธีทำงาน หรือการสื่อสาร การกำหนดข้อตกลงกลุ่มและทบทวนการทำงานเป็นระยะจึงเป็นแนวทางช่วยให้สมาชิกเข้าใจความรับผิดชอบร่วมกัน

แนวคิดดังกล่าวสอดคล้องกับฟีเจอร์ข้อตกลงกลุ่มของ Collab Reflection ซึ่งเปิดให้อาจารย์นำเอกสารข้อตกลงมาให้สมาชิกศึกษา และใช้การประเมินหลายรอบเพื่อสนับสนุนการสะท้อนพฤติกรรมระหว่างภาคการศึกษา ทั้งนี้ ระบบเป็นเครื่องมือช่วยติดตามและสื่อสาร ไม่ได้ทำหน้าที่แก้ไขความขัดแย้งแทนการพูดคุยระหว่างสมาชิก

\section{ปัญหาที่พบในการทำงานกลุ่ม}

\subsection{Social Loafing และ Free-rider}

Social Loafing หมายถึงปรากฏการณ์ที่บุคคลมีแนวโน้มใช้ความพยายามน้อยลงเมื่อทำงานเป็นกลุ่ม เมื่อเทียบกับการทำงานคนเดียว~\cite{latane1979many} หากสมาชิกทุกคนได้รับคะแนนจากผลงานกลุ่มเท่ากัน อาจเกิดสถานการณ์ที่สมาชิกบางคนมีส่วนร่วมน้อยแต่ได้รับผลลัพธ์เท่ากับสมาชิกที่ทำงานมากกว่า ซึ่งมักเรียกว่า Free-rider ขณะเดียวกันสมาชิกที่รับผิดชอบงานมากอาจต้องแบกรับภาระของกลุ่ม

ประเด็นนี้สอดคล้องกับผลการสัมภาษณ์อาจารย์ผู้สอน 5 ท่านในโครงงาน ซึ่งพบปัญหาสมาชิกไม่ช่วยงานและสมาชิกบางคนต้องรับภาระงานแทนผู้อื่น ดังนั้น Collab Reflection จึงให้สมาชิกประเมินพฤติกรรมการทำงานของตนเองและเพื่อนตามเกณฑ์เดียวกัน พร้อมแสดงผลรายด้านให้อาจารย์ใช้ประกอบการติดตาม โดยผลการประเมินไม่ควรถูกใช้เป็นข้อสรุปเกี่ยวกับบุคคลโดยปราศจากการพิจารณาบริบทอื่น

\subsection{ข้อจำกัดด้านอัตราส่วนผู้สอนต่อนักศึกษา}

ในรายวิชาที่มีนักศึกษาจำนวนมาก ผู้สอนอาจไม่สามารถสังเกตพฤติกรรมการทำงานของสมาชิกทุกคนได้อย่างทั่วถึง โดยเฉพาะการแบ่งงาน การสื่อสาร และการช่วยเหลือกันนอกชั้นเรียน ข้อมูลจากการสัมภาษณ์ในโครงงานสะท้อนว่าผู้สอนต้องการเครื่องมือที่เรียบง่าย ลดภาระการติดตาม และช่วยให้มองเห็นปัญหาของแต่ละทีมได้รวดเร็วขึ้น

Collab Reflection จึงออกแบบแดชบอร์ดสำหรับแสดงภาพรวมผลประเมิน ติดตามจำนวนผู้ส่ง และแจ้งสัญญาณจากรูปแบบคะแนน เช่น คะแนนเฉลี่ยที่ต่ำในบางด้านหรือความแตกต่างระหว่างคะแนนประเมินตนเองกับคะแนนจากเพื่อน สัญญาณดังกล่าวเป็นข้อมูลเพื่อช่วยจัดลำดับการติดตาม ไม่ใช่การวินิจฉัยหรือการตัดสินนักศึกษาโดยอัตโนมัติ

\section{การประเมินโดยเพื่อนร่วมกลุ่มและการประเมินตนเอง}

การประเมินโดยเพื่อน (Peer Assessment) เป็นกระบวนการที่ผู้เรียนพิจารณาคุณภาพหรือพฤติกรรมของเพื่อนในระดับเดียวกัน โดยอาศัยเกณฑ์หรือเป้าหมายที่กำหนดไว้ ส่วนการประเมินตนเอง (Self-assessment) เปิดโอกาสให้ผู้เรียนทบทวนผลงานและพฤติกรรมของตนเอง การใช้ทั้งสองวิธีร่วมกันช่วยให้ผู้เรียนได้รับมุมมองที่หลากหลายและมีส่วนรับผิดชอบต่อการเรียนรู้ของตนเอง~\cite{topping1998peer}

ในบริบทของโครงงานนี้ การประเมินเพื่อนไม่ได้มุ่งเป็นเครื่องมือตัดเกรดเพียงอย่างเดียว แต่เป็นพื้นที่สะท้อนพฤติกรรมและการมีส่วนร่วมของสมาชิก โดยเน้นข้อเสนอแนะเชิงสร้างสรรค์ (Constructive Feedback) ที่ผู้รับสามารถนำไปปรับปรุงการทำงานในรอบถัดไปได้ อย่างไรก็ตาม การประเมินอาจได้รับอิทธิพลจากความสนิทสนม ความเกรงใจ หรือความขัดแย้งส่วนบุคคล ระบบจึงออกแบบให้ปกปิดตัวตนผู้ประเมิน ใช้ข้อมูลจากหลายมุมมอง และให้อาจารย์เป็นผู้ตรวจทานก่อนเผยแพร่ผล

\subsection{การประเมินระหว่างทางและการประเมินเพื่อสรุปผล}\label{sec:ch2-formative}

การประเมินระหว่างทาง (Formative Assessment) มีเป้าหมายเพื่อให้ข้อเสนอแนะระหว่างกระบวนการเรียนรู้ เพื่อนำไปปรับปรุงการทำงาน ส่วนการประเมินเพื่อสรุปผล (Summative Assessment) ใช้พิจารณาผลสัมฤทธิ์เมื่อสิ้นสุดช่วงการเรียนรู้~\cite{black1998assessment} แนวคิดนี้เป็นพื้นฐานของการกำหนดประเภทรอบประเมินใน Collab Reflection โดยอาจารย์สามารถกำหนดรอบฝึกเพื่อการพัฒนา และรอบเก็บคะแนนตามวัตถุประสงค์ของรายวิชาได้

\subsection{ระบบประเมินโดยเพื่อนที่เกี่ยวข้อง}

ระบบที่มีอยู่เดิมเป็นแนวทางอ้างอิงในการกำหนดความสามารถของ Collab Reflection โดยแต่ละระบบมีจุดเน้นแตกต่างกัน ดังตารางที่~\ref{tab:ch2-1}

\begin{table}[htbp]
\centering
\caption{การเปรียบเทียบระบบประเมินโดยเพื่อนร่วมกลุ่มที่เกี่ยวข้อง}\label{tab:ch2-1}
\begin{tabularx}{\textwidth}{P{0.2\textwidth}YY}
\hline
\textbf{ระบบ} & \textbf{ลักษณะการทำงานตามเอกสารที่ศึกษา} & \textbf{ประเด็นที่นำมาใช้เป็นแนวทาง} \\ \hline
CATME Peer Evaluation & ประเมินสมาชิกในทีมตามพฤติกรรมหลายด้าน รองรับการประเมินตนเอง การประเมินเพื่อน และความคิดเห็นเป็นข้อความ~\cite{ohland2012catme} & ใช้กรอบพฤติกรรม 5 ด้านเป็นพื้นฐานของเกณฑ์ประเมิน \\
FeedbackFruits Peer Review & รองรับการประเมินผลงานตามเกณฑ์และ Rubric การรีวิวแบบไม่เปิดเผยตัวตน และมี AI Feedback Coach & ใช้เป็นตัวอย่างการผสานเกณฑ์ประเมินกับเครื่องมือสนับสนุนการให้ข้อเสนอแนะ \\
Peerceptiv Peer Assessment & รองรับการให้คะแนนและเขียนความคิดเห็นเกี่ยวกับผลงาน รวมถึงการประเมินการมีส่วนร่วม การสื่อสาร และการทำงานร่วมกัน ตลอดจนการรีวิวแบบไม่เปิดเผยตัวตน & ใช้เป็นตัวอย่างระบบที่รวบรวมการประเมินจากผู้เรียนหลายคน \\ \hline
\end{tabularx}
\end{table}

จากแนวทางข้างต้น Collab Reflection มุ่งรวมการประเมินตนเอง การประเมินเพื่อนในทีม เกณฑ์ที่อธิบายพฤติกรรม การสะท้อนคิดแบบคำถามปลายเปิด และแดชบอร์ดสำหรับอาจารย์ไว้ในแพลตฟอร์มเดียว ส่วนการใช้ปัญญาประดิษฐ์ (AI) เพื่อช่วยตรวจโทนข้อความและสรุปความคิดเห็นถูกกำหนดเป็นฟีเจอร์ในระยะที่ 2 ตามแผนการพัฒนาของระบบ

\section{การสะท้อนคิดและการให้ข้อเสนอแนะ}

การสะท้อนคิด (Reflection) เป็นกระบวนการทบทวนประสบการณ์เพื่อสร้างความเข้าใจและนำสิ่งที่เรียนรู้ไปปรับใช้ในการปฏิบัติครั้งต่อไป Kolb อธิบายวงจรการเรียนรู้จากประสบการณ์ไว้ 4 ขั้น ได้แก่ ประสบการณ์รูปธรรม การสังเกตและสะท้อนคิด การสร้างแนวคิดนามธรรม และการทดลองปฏิบัติ~\cite{kolb1984experiential} ส่วน Schön แบ่งการสะท้อนคิดเป็นการสะท้อนคิดขณะปฏิบัติ (Reflection-in-action) และการสะท้อนคิดหลังปฏิบัติ (Reflection-on-action)~\cite{schon1983reflective}

Collab Reflection เน้นการสะท้อนคิดหลังปฏิบัติ โดยให้นักศึกษาทบทวนการทำงานของตนเองและเพื่อนเป็นรอบตามช่วงเวลาที่อาจารย์กำหนด คำถามปลายเปิดใช้แนวคิด ``ให้ดอกไม้ เติมปุ๋ย'' โดย ``ดอกไม้'' ให้ระบุสิ่งที่สมาชิกทำได้ดี ส่วน ``ปุ๋ย'' ให้เสนอสิ่งที่ช่วยให้สมาชิกทำได้ดีขึ้นในงานถัดไป สำหรับการประเมินตนเองจะปรับเป็นการระบุสิ่งที่ตนทำได้ดีและสิ่งที่ต้องการปรับปรุง

Hattie และ Timperley อธิบายว่าข้อเสนอแนะที่มีประสิทธิภาพควรช่วยให้ผู้เรียนเข้าใจเป้าหมาย ความก้าวหน้าในปัจจุบัน และแนวทางดำเนินการต่อไป~\cite{hattie2007feedback} ดังนั้นระบบจึงแนะนำให้ผู้ประเมินกล่าวถึงพฤติกรรมหรือผลงานที่สังเกตได้จริง และเสนอแนวทางที่นำไปปฏิบัติได้ แทนการตำหนิหรือตัดสินตัวบุคคล

\section{เกณฑ์การประเมินแบบ CATME}

โครงงานเลือกใช้กรอบการประเมินสมาชิกทีมของ CATME (Comprehensive Assessment of Team Member Effectiveness) เป็นพื้นฐาน โดย CATME พัฒนาเกณฑ์แบบมาตราส่วนที่มีคำอธิบายพฤติกรรมกำกับ (Behaviorally Anchored Rating Scale: BARS) สำหรับการประเมินตนเองและเพื่อน~\cite{ohland2012catme} กรอบดังกล่าวแบ่งพฤติกรรมการทำงานเป็นทีมออกเป็น 5 ด้าน ดังตารางที่~\ref{tab:ch2-2}

\begin{table}[htbp]
\centering
\caption{องค์ประกอบการประเมินพฤติกรรมการทำงานเป็นทีมตามแนวคิด CATME}\label{tab:ch2-2}
\begin{tabularx}{\textwidth}{P{0.32\textwidth}Yl}
\hline
\textbf{ด้านการประเมิน} & \textbf{พฤติกรรมที่พิจารณา} & \textbf{กลุ่มด้าน} \\ \hline
การมีส่วนร่วมในงาน (Contributing to the Team's Work) & ทำงานที่ได้รับมอบหมายให้ครบถ้วนและตรงเวลา รวมถึงช่วยเหลือเพื่อนเมื่อจำเป็น & ด้านงาน \\
การปฏิสัมพันธ์กับเพื่อน (Interacting with Teammates) & รับฟังความคิดเห็น สื่อสารอย่างชัดเจน และปฏิบัติต่อสมาชิกอย่างเหมาะสม & ด้านทีม \\
การดูแลให้ทีมดำเนินตามแผน (Keeping the Team on Track) & ติดตามความคืบหน้า ประสานงาน และช่วยให้ทีมดำเนินงานตามแผน & ด้านทีม \\
การคาดหวังคุณภาพ (Expecting Quality) & ใส่ใจคุณภาพของงานและสนับสนุนให้ผลงานของทีมมีคุณภาพ & ด้านงาน \\
ความรู้และทักษะที่เกี่ยวข้อง (Having Relevant KSAs) & มีความรู้ ทักษะ และความสามารถที่จำเป็นต่อการทำงานของทีม & ด้านงาน \\ \hline
\end{tabularx}
\end{table}

ในระบบกำหนดคะแนนระดับ 1--5 โดยมีคำอธิบายพฤติกรรมที่ระดับ 1, 3 และ 5 ส่วนระดับ 2 และ 4 เป็นค่าระหว่างกลาง เพื่อให้ผู้ประเมินใช้พฤติกรรมที่สังเกตได้เป็นหลักในการเลือกคะแนน ตัวอย่างเช่น ระดับ 1 หมายถึงพฤติกรรมต่ำกว่าที่คาดหวัง ระดับ 3 หมายถึงทำได้ตามที่คาดหวัง และระดับ 5 หมายถึงทำได้เกินกว่าที่คาดหวัง

นอกจากคะแนนทั้ง 5 ด้านแล้ว แบบประเมินยังมีคำถามปลายเปิดเพื่อระบุจุดแข็งและสิ่งที่ควรพัฒนา ระบบคำนวณค่าเฉลี่ยคะแนนที่เพื่อนให้ในแต่ละด้านโดยไม่รวมคะแนนประเมินตนเอง แล้วนำมาเปรียบเทียบกับคะแนนที่นักศึกษาประเมินตนเอง เพื่อให้เห็นความแตกต่างระหว่างมุมมอง ทั้งนี้การเปรียบเทียบมีไว้เพื่อการสะท้อนคิด ไม่ใช่ข้อสรุปว่ามุมมองใดถูกต้องกว่า

การคำนวณคะแนนถ่วงน้ำหนักร่วมกับคะแนนชิ้นงานเป็นฟีเจอร์ที่วางแผนไว้ในระยะที่ 3 ตามการออกแบบในบทที่~\ref{ch:method} จึงไม่ถือเป็นความสามารถของระบบในระยะที่ 1

\section{การจัดการข้อเสนอแนะที่อ่อนไหวและการคุ้มครองข้อมูลส่วนบุคคล}

การประเมินเพื่อนเกี่ยวข้องกับความคิดเห็นต่อพฤติกรรมของบุคคล จึงอาจมีข้อความที่รุนแรง ตำหนิตัวบุคคล หรือสะท้อนความขัดแย้งภายในกลุ่ม การออกแบบระบบจึงควรลดโอกาสที่ข้อเสนอแนะจะสร้างผลกระทบต่อความสัมพันธ์ระหว่างสมาชิก และสนับสนุนให้ผู้ใช้กล่าวถึงพฤติกรรมหรือผลงานที่สังเกตได้ แทนการกล่าวโจมตีตัวบุคคล

ในระยะที่ 1 ระบบใช้คำแนะนำและตัวอย่างการเขียนบนแบบฟอร์ม รวมถึงแนวคิด ``ให้ดอกไม้ เติมปุ๋ย'' เพื่อช่วยให้ผู้ประเมินเขียนข้อเสนอแนะอย่างสร้างสรรค์ ส่วนปัญญาประดิษฐ์ (AI) สำหรับตรวจโทนข้อความและเสนอถ้อยคำทางเลือกเป็นฟีเจอร์ในระยะที่ 2 โดยปัญญาประดิษฐ์ (AI) ทำหน้าที่แจ้งเตือนและเสนอทางเลือกเท่านั้น ผู้เขียนเป็นผู้ตัดสินใจว่าจะใช้ข้อความที่เสนอหรือไม่ และอาจารย์ยังเป็นผู้ตรวจทานก่อนเผยแพร่ผล

เนื่องจากข้อมูลการประเมินอาจเชื่อมโยงกับบุคคลได้ การเก็บ ใช้ และเปิดเผยข้อมูลจึงต้องคำนึงถึงพระราชบัญญัติคุ้มครองข้อมูลส่วนบุคคล พ.ศ. 2562~\cite{pdpa2562} ระบบกำหนดให้แจ้งวัตถุประสงค์ ประเภทข้อมูล ผู้ที่สามารถเข้าถึงข้อมูล ระยะเวลาการเก็บรักษา และสิทธิของเจ้าของข้อมูล พร้อมเก็บข้อมูลเท่าที่จำเป็นต่อการประเมินการทำงานกลุ่ม

มาตรการที่ออกแบบไว้ประกอบด้วยการเข้าสู่ระบบด้วยบัญชีมหาวิทยาลัย การกำหนดสิทธิ์ตามบทบาท การปกปิดตัวตนผู้ประเมินจากเพื่อน และการซ่อนผลประเมินเมื่อมีผู้ประเมินไม่ถึง 2 คน เพื่อลดความเสี่ยงในการคาดเดาตัวตน นอกจากนี้ นักศึกษาจะเห็นเฉพาะข้อมูลของตนเอง ขณะที่อาจารย์สามารถตรวจสอบผลในขอบเขตของรายวิชาที่รับผิดชอบ และเป็นผู้ตัดสินใจเผยแพร่ผล

สำหรับฟีเจอร์ปัญญาประดิษฐ์ (AI) ในระยะที่ 2 ระบบยึดหลัก Human-in-the-loop โดยไม่ส่งชื่อหรือรหัสนักศึกษาไปยังผู้ให้บริการโมเดล และส่งเฉพาะข้อความที่จำเป็นต่อการประมวลผล ทั้งนี้ ผู้ใช้สามารถเลือกยินยอมในส่วนทางเลือกตามการออกแบบความยินยอมของระบบ และอาจารย์ยังคงเป็นผู้ตรวจทานข้อเสนอแนะก่อนส่งทุกครั้ง

\section{แนวคิดด้านสถาปัตยกรรมและเทคโนโลยีที่เกี่ยวข้อง}

Collab Reflection พัฒนาในรูปแบบเว็บแอปพลิเคชันที่แยกส่วนหน้า (Frontend) และส่วนหลัง (Backend) ออกจากกัน โดยสื่อสารผ่าน REST API และใช้ JSON เป็นรูปแบบข้อมูล การแยกส่วนดังกล่าวช่วยให้แต่ละส่วนพัฒนาและทดสอบได้เป็นอิสระ รวมถึงกำหนดขอบเขตความรับผิดชอบระหว่างส่วนติดต่อผู้ใช้กับตรรกะการทำงานของระบบได้ชัดเจน

ส่วนหน้าพัฒนาด้วย React 19, TypeScript และ Vite ใช้ Tailwind CSS ในการจัดรูปแบบหน้าจอ ใช้ TanStack Query จัดการข้อมูลจากเซิร์ฟเวอร์ และ Zustand จัดการสถานะภายในแอปพลิเคชัน ส่วนหลังพัฒนาด้วย Node.js และ Express 5 ให้บริการ API และใช้ Drizzle ORM เชื่อมต่อกับฐานข้อมูล PostgreSQL ซึ่งจัดเก็บข้อมูลผู้ใช้ รายวิชา กลุ่ม รอบประเมิน และผลการประเมิน

การติดตั้งระบบใช้ Docker เพื่อจัดการสภาพแวดล้อมของบริการ และ nginx สำหรับให้บริการไฟล์ส่วนหน้า ทำหน้าที่ reverse proxy และจัดการ HTTPS การทดสอบประกอบด้วย Vitest สำหรับการทดสอบระดับหน่วยและการเชื่อมต่อ และ Cypress สำหรับการทดสอบการทำงานตั้งแต่ต้นจนจบ

ในระยะที่ 2 มีแผนเพิ่ม worker สำหรับประมวลผลงานปัญญาประดิษฐ์ (AI) ผ่านคิวงานเบื้องหลัง โดยใช้โมเดลภาษา Qwen ผ่าน API และบันทึกผลกลับฐานข้อมูล การใช้คิวช่วยไม่ให้ผู้ใช้ต้องรอการประมวลผลในหน้าจอโดยตรง ทั้งนี้ ระบบเลือกใช้ PostgreSQL สำหรับจัดเก็บข้อมูลและรองรับคิวงาน เพื่อลดจำนวนบริการที่ต้องดูแลให้เหมาะกับทรัพยากรเซิร์ฟเวอร์ของโครงงาน

ด้านการเชื่อมต่อระบบภายนอก โครงงานวางแผนรองรับการเชื่อมต่อ Mango Canvas ผ่านมาตรฐาน Learning Tools Interoperability (LTI) ในระยะที่ 3 ส่วนระยะที่ 1 ใช้การนำเข้ารายชื่อนักศึกษาจากไฟล์ CSV เพื่อให้สามารถเริ่มใช้งานได้โดยไม่ต้องรอการตั้งค่าจากผู้ดูแลระบบมหาวิทยาลัย

\section{ช่องว่างของระบบที่มีอยู่และแนวทางของ Collab Reflection}

จากการศึกษาแนวคิดและระบบที่เกี่ยวข้อง พบว่าระบบประเมินเพื่อนที่มีอยู่มีความสามารถด้านการให้คะแนน การใช้ Rubric การแสดงความคิดเห็น หรือการประเมินการมีส่วนร่วมในทีมในรูปแบบที่แตกต่างกัน อย่างไรก็ตาม โครงงานนี้มุ่งตอบโจทย์เฉพาะบริบทของรายวิชาที่มีนักศึกษาจำนวนมากและต้องการติดตามพฤติกรรมการทำงานเป็นทีมอย่างต่อเนื่อง โดยผสานการประเมินตนเองและเพื่อนในทีมเข้ากับคำถามสะท้อนคิดและเครื่องมือช่วยให้อาจารย์เห็นภาพรวม

แนวทางของ Collab Reflection จึงประกอบด้วย (1) เกณฑ์ประเมินพฤติกรรมการทำงานเป็นทีม 5 ด้านที่มีคำอธิบายคะแนน (2) การประเมินตนเองและเพื่อนร่วมทีมในรอบที่อาจารย์กำหนด (3) คำถามปลายเปิดเพื่อชื่นชมจุดแข็งและเสนอแนวทางพัฒนา (4) การปกปิดตัวตนและกำหนดสิทธิ์การเข้าถึงข้อมูล (5) แดชบอร์ดและสัญญาณเตือนเพื่อช่วยอาจารย์ติดตามทีม และ (6) การวางแผนเพิ่มปัญญาประดิษฐ์ (AI) เพื่อช่วยตรวจโทนและสรุปความคิดเห็น โดยคงให้อาจารย์เป็นผู้ตรวจทานขั้นสุดท้าย

ดังนั้น จุดเน้นของโครงงานไม่ได้อยู่ที่การสร้างระบบให้คะแนนเพียงอย่างเดียว แต่เป็นการออกแบบกระบวนการสะท้อนคิดที่นำไปใช้พัฒนาทักษะการทำงานเป็นทีมได้ พร้อมลดภาระการรวบรวมและสรุปข้อมูลของผู้สอน ทั้งนี้ การเชื่อมต่อกับระบบมหาวิทยาลัย การคำนวณคะแนนถ่วงน้ำหนัก และการติดตามพัฒนาการสะสมยังเป็นแผนพัฒนาระยะถัดไปตามที่ระบุในบทที่~\ref{ch:method}

\section{ความรู้ตามหลักสูตรที่นำมาใช้หรือบูรณาการในโครงงาน}

ความรู้ตามหลักสูตรวิศวกรรมคอมพิวเตอร์ที่นำมาใช้ประกอบด้วย

\begin{itemize}
  \item \textbf{การพัฒนาเว็บแอปพลิเคชัน:} การออกแบบส่วนติดต่อผู้ใช้ การพัฒนาส่วนหน้าและส่วนหลัง และการเชื่อมต่อผ่าน REST API
  \item \textbf{ระบบฐานข้อมูล:} การออกแบบฐานข้อมูลเชิงสัมพันธ์ การกำหนดความสัมพันธ์ระหว่างข้อมูล และการใช้ SQL เพื่อจัดเก็บและสรุปผลการประเมิน
  \item \textbf{วิศวกรรมซอฟต์แวร์:} การเก็บและวิเคราะห์ความต้องการ การกำหนดขอบเขตและลำดับการพัฒนา การออกแบบกระบวนการใช้งาน และการทดสอบระบบ
  \item \textbf{การปฏิสัมพันธ์ระหว่างมนุษย์กับคอมพิวเตอร์ (HCI):} การออกแบบ UX/UI ให้เหมาะกับผู้ใช้ 2 กลุ่ม ได้แก่ นักศึกษาและอาจารย์
  \item \textbf{การเขียนโปรแกรมและเครือข่ายคอมพิวเตอร์:} การพัฒนาโปรแกรมฝั่งเซิร์ฟเวอร์ และการสื่อสารระหว่าง Client กับ Server ผ่าน HTTP
\end{itemize}

\section{ความรู้นอกหลักสูตรที่นำมาใช้หรือบูรณาการในโครงงาน}

\begin{itemize}
  \item \textbf{เครื่องมือพัฒนาส่วนหน้าเว็บ:} React 19, Vite, TypeScript, Tailwind CSS, TanStack Query และ Zustand
  \item \textbf{เครื่องมือพัฒนาส่วนหลังบ้าน:} Node.js, Express 5 และ Drizzle ORM สำหรับเชื่อมต่อกับ PostgreSQL
  \item \textbf{โมเดลภาษาขนาดใหญ่ (LLM):} การเรียกใช้ Qwen API เพื่อช่วยตรวจโทนข้อความ เสนอถ้อยคำเชิงสร้างสรรค์ และสรุปความคิดเห็นตามแนวทาง Human-in-the-loop ในระยะที่ 2
  \item \textbf{DevOps และการทดสอบ:} การใช้ Docker และ nginx ในการติดตั้งระบบ การทดสอบด้วย Vitest และการทดสอบ End-to-End ด้วย Cypress
  \item \textbf{การยืนยันตัวตนและการเชื่อมต่อระบบ:} การเข้าสู่ระบบด้วยบัญชี CMU ผ่าน CMU OAuth/\allowbreak OpenID Connect และแผนการเชื่อมต่อ Mango Canvas ผ่าน LTI ในระยะที่ 3
  \item \textbf{มาตรฐานและแนวปฏิบัติ:} REST API, JSON, OpenAPI Specification, SQL และ WCAG 2.2 เพื่อสนับสนุนการเชื่อมต่อ การจัดการข้อมูล และการเข้าถึงระบบ
  \item \textbf{การคุ้มครองข้อมูลส่วนบุคคล:} การประยุกต์ใช้หลักการเก็บข้อมูลเท่าที่จำเป็น การแจ้งรายละเอียดการใช้ข้อมูล การกำหนดสิทธิ์เข้าถึง และการจัดการความยินยอมตาม PDPA
  \item \textbf{การออกแบบโดยยึดผู้ใช้เป็นศูนย์กลาง:} การสัมภาษณ์อาจารย์ผู้สอน 5 ท่าน วิเคราะห์ปัญหาและความต้องการ แล้วนำผลมาใช้กำหนดฟีเจอร์และต้นแบบหน้าจอ
\end{itemize}
